\chapter{Cơ sở lý thuyết}
\label{ch:background}

Chương này trình bày các kiến thức nền tảng cần thiết để hiểu phương pháp đề xuất ở
Chương~\ref{ch:method}. Trước hết, chúng tôi giới thiệu đồ thị tri thức và mở rộng có yếu tố thời
gian của nó, cùng phát biểu hình thức cho bài toán hỏi đáp trên đồ thị tri thức thời gian. Tiếp
theo, chúng tôi tóm lược về mô hình ngôn ngữ lớn và các kỹ thuật tăng cường năng lực suy luận cho
nó, gồm tăng cường tri thức, suy luận có cấu trúc và cơ chế bộ nhớ -- nền tảng triết học mà phương
pháp Suy luận Trừu tượng Quy nạp kế thừa. Cuối cùng, chương trình bày một số công cụ kỹ thuật được
sử dụng trực tiếp trong khóa luận: phân cụm K-means và biểu diễn ngữ nghĩa, kiến trúc cross-encoder
và bài toán nhận dạng thực thể có tên cho tiếng Việt. Bảng~\ref{tab:notation} tổng hợp các ký hiệu
toán học dùng xuyên suốt khóa luận.

\section{Đồ thị tri thức và đồ thị tri thức thời gian}
\label{sec:tkg}

\subsection{Đồ thị tri thức}

Đồ thị tri thức (\textit{Knowledge Graph} -- KG) là một cách biểu diễn tri thức dưới dạng đồ thị,
trong đó các nút biểu diễn thực thể và các cạnh biểu diễn quan hệ giữa chúng. Tri thức được mã hóa
thành tập các \textit{bộ ba} (\textit{triple}) có dạng (chủ thể, quan hệ, đối tượng).

\begin{definition}[Đồ thị tri thức]
Một đồ thị tri thức là một bộ $\mathcal{G} = (\mathcal{E}, \mathcal{R}, \mathcal{F})$, trong đó
$\mathcal{E}$ là tập thực thể, $\mathcal{R}$ là tập quan hệ, và
$\mathcal{F} \subseteq \mathcal{E} \times \mathcal{R} \times \mathcal{E}$ là tập sự kiện. Mỗi sự
kiện là một bộ ba $(s, r, o)$ với $s, o \in \mathcal{E}$ lần lượt là chủ thể và đối tượng, còn
$r \in \mathcal{R}$ là quan hệ nối giữa chúng.
\end{definition}

Ví dụ, bộ ba (\textit{Barack Obama}, \texttt{chức\_vụ}, \textit{Tổng thống Hoa Kỳ}) khẳng định
rằng Barack Obama giữ chức Tổng thống Hoa Kỳ. Các đồ thị tri thức quy mô lớn như
Wikidata~\cite{vrandecic2014wikidata}, DBpedia hay YAGO đã trở thành nguồn tri thức có cấu trúc
quan trọng cho nhiều ứng dụng xử lý ngôn ngữ tự nhiên, đặc biệt là hỏi đáp.

\subsection{Đồ thị tri thức thời gian}

Hạn chế của đồ thị tri thức tĩnh là không biểu diễn được khía cạnh thời gian của tri thức. Trong
thực tế, phần lớn sự kiện chỉ đúng trong một khoảng thời gian nhất định: một người chỉ giữ một
chức vụ trong một nhiệm kỳ, hai quốc gia chỉ ở trạng thái xung đột trong một giai đoạn lịch sử.
\textit{Đồ thị tri thức thời gian} (\textit{Temporal Knowledge Graph} -- TKG) mở rộng đồ thị tri
thức bằng cách gắn thêm thông tin thời gian cho mỗi sự kiện.

\begin{definition}[Đồ thị tri thức thời gian]
Một đồ thị tri thức thời gian là một bộ $\mathcal{K} = (\mathcal{E}, \mathcal{R}, \mathcal{T},
\mathcal{F})$, trong đó $\mathcal{E}$, $\mathcal{R}$ lần lượt là tập thực thể và quan hệ,
$\mathcal{T}$ là tập mốc thời gian, và $\mathcal{F}$ là tập sự kiện thời gian. Mỗi sự kiện là một
bộ bốn $(s, r, o, \tau)$ với $s, o \in \mathcal{E}$, $r \in \mathcal{R}$ và $\tau \in \mathcal{T}$
là mốc (hoặc khoảng) thời gian mà sự kiện có hiệu lực.
\end{definition}

Khi $\tau$ là một khoảng, ta viết $\tau = [t_{b}, t_{e}]$ với $t_{b}$ là thời điểm bắt đầu và
$t_{e}$ là thời điểm kết thúc. Chẳng hạn, sự kiện (\textit{Barack Obama}, \texttt{chức\_vụ},
\textit{Tổng thống Hoa Kỳ}, $[2009, 2017]$) cho biết khoảng thời gian Obama tại nhiệm.
Hình~\ref{fig:tkg-example} minh họa một đồ thị tri thức thời gian nhỏ với nhiều quan hệ gắn mốc
thời gian khác nhau.

\begin{figure}[ht]
  \centering
  \fbox{\parbox{0.85\textwidth}{\centering\vspace{2mm}\footnotesize
  \textit{G.~W.~Bush} $\xrightarrow{\text{tổng thống Hoa Kỳ}~[2001,2009]}$ \textit{Hoa Kỳ}\\[1mm]
  \textit{B.~Obama} $\xrightarrow{\text{tổng thống Hoa Kỳ}~[2009,2017]}$ \textit{Hoa Kỳ}\\[1mm]
  \textit{Hoa Kỳ} $\xrightarrow{\text{thành viên}~[1945,-]}$ \textit{Liên Hợp Quốc}
  \vspace{2mm}}}
  \caption{Một đồ thị tri thức thời gian thu nhỏ: cùng một quan hệ
  \texttt{tổng\_thống} nối nhiều chủ thể tới cùng đối tượng nhưng ở những khoảng thời gian khác
  nhau. \todo{thay bằng hình đồ thị trực quan có nút và cạnh}}
  \label{fig:tkg-example}
\end{figure}

\subsection{Biểu diễn thời gian trong Wikidata}

Trong khóa luận, bộ dữ liệu được trích xuất từ Wikidata. Ở Wikidata, mỗi thực thể được định danh
bởi một mã \texttt{QID} (ví dụ \texttt{Q76} cho Barack Obama) và mỗi thuộc tính (quan hệ) bởi một
mã \texttt{PID} (ví dụ \texttt{P39} cho ``chức vụ nắm giữ''). Khía cạnh thời gian được biểu diễn
thông qua các \textit{định tố} (\textit{qualifier}) gắn vào mỗi khẳng định, trong đó ba định tố
quan trọng nhất là \texttt{P580} (thời điểm bắt đầu), \texttt{P582} (thời điểm kết thúc) và
\texttt{P585} (thời điểm diễn ra). Việc khai thác các định tố này cho phép tái dựng các sự kiện
thời gian dưới dạng bộ bốn $(s, r, o, \tau)$ phục vụ cho bài toán TKGQA, sẽ được trình bày chi
tiết ở Mục~\ref{sec:cronqvn}.

\section{Bài toán hỏi đáp trên đồ thị tri thức thời gian}
\label{sec:tkgqa-def}

Hỏi đáp trên đồ thị tri thức thời gian (\textit{Temporal Knowledge Graph Question Answering} --
TKGQA) là bài toán trả lời câu hỏi ngôn ngữ tự nhiên có ràng buộc về thời gian, dựa trên một đồ
thị tri thức thời gian cho trước.

\begin{definition}[Bài toán TKGQA]
Cho một đồ thị tri thức thời gian $\mathcal{K} = (\mathcal{E}, \mathcal{R}, \mathcal{T},
\mathcal{F})$ và một câu hỏi ngôn ngữ tự nhiên $q$, bài toán TKGQA yêu cầu trả về một đáp án
$a$ là một thực thể $a \in \mathcal{E}$ hoặc một mốc thời gian $a \in \mathcal{T}$ sao cho $a$
trả lời đúng câu hỏi $q$ dựa trên tri thức chứa trong $\mathcal{K}$.
\end{definition}

Ví dụ, với câu hỏi $q = $ \textit{``Ai giữ chức Tổng thống Hoa Kỳ vào năm 2010?''}, đáp án là
thực thể \textit{Barack Obama}; còn với câu hỏi \textit{``Barack Obama bắt đầu giữ chức Tổng thống
Hoa Kỳ năm nào?''}, đáp án là mốc thời gian \textit{2009}.

Các câu hỏi TKGQA thường được phân loại theo hai chiều~\cite{saxena2021cronkgqa, chen2023multitq}.
Theo \textbf{độ phức tạp suy luận}, câu hỏi được chia thành \textit{đơn giản} (\textit{simple})
-- chỉ cần một bước truy vấn trực tiếp trên một sự kiện -- và \textit{phức tạp}
(\textit{complex}) -- cần kết hợp nhiều sự kiện, so sánh hoặc lọc theo thời gian. Theo
\textbf{kiểu đáp án}, câu hỏi được chia thành loại trả về \textit{thực thể} (\textit{entity}) và
loại trả về \textit{thời gian} (\textit{time}). Bộ dữ liệu CronQ-VN được xây dựng trong khóa luận
tổ chức câu hỏi thành năm loại cụ thể, sẽ được trình bày ở Mục~\ref{sec:generate}.

Điểm mấu chốt khiến TKGQA khó hơn hỏi đáp thông thường là yêu cầu \textit{suy luận nhiều bước có
ràng buộc thời gian}: hệ thống không chỉ truy xuất một sự kiện mà phải xác định mốc thời gian
tham chiếu, rồi lọc và so sánh các sự kiện theo trình tự thời gian. Đây chính là điểm yếu của mô
hình ngôn ngữ lớn khi dùng trực tiếp, như sẽ phân tích ở các mục tiếp theo.

\subsection{Phân loại câu hỏi và thao tác suy luận}

Bảng~\ref{tab:qtypes} hệ thống hóa các loại câu hỏi thời gian thường gặp cùng thao tác suy luận
cốt lõi tương ứng. Cách phân loại này vừa định hướng việc sinh dữ liệu (Mục~\ref{sec:generate}),
vừa là cơ sở để phân tích kết quả chi tiết theo từng loại ở Chương~\ref{ch:experiment}. Có thể
thấy các loại câu hỏi đòi hỏi mức độ suy luận tăng dần: từ truy vấn trực tiếp một sự kiện, đến lọc
theo mốc tham chiếu, sắp xếp theo thời gian, và phức tạp nhất là giao nhiều ràng buộc thời gian.

\begin{table}[ht]
\centering
\caption{Các loại câu hỏi thời gian và thao tác suy luận cốt lõi.}
\label{tab:qtypes}
\begin{tabular}{L{3.2cm} L{5.3cm} L{4.0cm}}
\toprule
\textbf{Loại câu hỏi} & \textbf{Mô tả} & \textbf{Thao tác chính} \\
\midrule
Thực thể đơn & Ai/cái gì giữ vai trò $r$ vào thời điểm $\tau$? & Truy vấn trực tiếp một bước \\
Thời gian đơn & $s$ giữ vai trò $r$ vào thời điểm nào? & Truy vấn mốc thời gian \\
Trước/Sau & Ai là $r$ trước/sau khi $s$ là $r$? & Lấy mốc tham chiếu, lọc thời gian \\
Đầu/Cuối & $s$ đầu tiên/gần nhất là gì? & Sắp xếp theo thời gian \\
Đồng thời & Ai là $r_1$ khi $s$ đang là $r_2$? & Giao theo khoảng thời gian \\
\bottomrule
\end{tabular}
\end{table}

\subsection{Độ đo đánh giá}
\label{sec:metrics}

Do đáp án của bài toán TKGQA là một thực thể hoặc một mốc thời gian xác định, độ đo phổ biến nhất
là \textit{độ chính xác} (\textit{accuracy}) -- tỉ lệ câu hỏi được trả lời đúng:

\begin{equation}
\label{eq:accuracy}
\mathrm{Acc} = \frac{1}{|Q|} \sum_{q \in Q} \mathbf{1}\!\left[\hat{a}_q \in \mathcal{A}_q\right],
\end{equation}

trong đó $Q$ là tập câu hỏi kiểm thử, $\hat{a}_q$ là đáp án dự đoán, $\mathcal{A}_q$ là tập đáp án
đúng (một câu hỏi có thể có nhiều đáp án đúng), và $\mathbf{1}[\cdot]$ là hàm chỉ thị nhận giá trị
$1$ khi điều kiện đúng và $0$ khi sai. Khi mô hình trả về một danh sách đáp án có xếp hạng, người
ta còn dùng độ đo \textit{Hits@1} -- tỉ lệ câu hỏi mà đáp án đúng nằm ở vị trí xếp hạng cao nhất.
Trong khóa luận, kết quả được báo cáo bằng độ chính xác tổng thể, đồng thời tách nhỏ theo từng loại
câu hỏi và theo kiểu đáp án (thực thể so với thời gian) nhằm làm rõ điểm mạnh, điểm yếu của phương
pháp trên từng dạng suy luận.

\section{Mô hình ngôn ngữ lớn và học trong ngữ cảnh}
\label{sec:llm}

\subsection{Kiến trúc Transformer}

Mô hình ngôn ngữ lớn (\textit{Large Language Model} -- LLM) hiện đại hầu hết dựa trên kiến trúc
Transformer~\cite{vaswani2017attention}, với thành phần cốt lõi là cơ chế \textit{tự chú ý}
(\textit{self-attention}). Cho một chuỗi biểu diễn đầu vào, cơ chế chú ý tính toán đầu ra theo
công thức:

\begin{equation}
\label{eq:attention}
\mathrm{Attention}(Q, K, V) = \mathrm{softmax}\!\left(\frac{Q K^{\top}}{\sqrt{d_k}}\right) V,
\end{equation}

trong đó $Q$, $K$, $V$ lần lượt là các ma trận truy vấn (\textit{query}), khóa (\textit{key}) và
giá trị (\textit{value}) được chiếu tuyến tính từ biểu diễn đầu vào, còn $d_k$ là số chiều của
vector khóa. Cơ chế này cho phép mỗi vị trí trong chuỗi tổng hợp thông tin từ mọi vị trí khác,
nắm bắt được các phụ thuộc tầm xa trong văn bản.

\subsection{Sinh văn bản tự hồi quy}

Mô hình ngôn ngữ lớn được huấn luyện theo lối \textit{tự hồi quy} (\textit{autoregressive}): với
chuỗi token $x_1, x_2, \dots, x_n$, mô hình ước lượng xác suất của chuỗi bằng tích các xác suất
có điều kiện,

\begin{equation}
\label{eq:autoregressive}
P(x_1, \dots, x_n) = \prod_{i=1}^{n} P(x_i \mid x_1, \dots, x_{i-1}).
\end{equation}

Tại mỗi bước sinh, mô hình chọn token tiếp theo dựa trên phân phối xác suất có điều kiện. Chính
cơ chế sinh \textit{tham lam theo xác suất cực đại} này khiến mô hình dễ tích lũy lỗi qua nhiều
bước khi phải thực hiện chuỗi suy luận dài, như đã nêu ở Mục~\ref{sec:limitations}.

\subsection{Chiến lược giải mã}

Từ phân phối xác suất ở Công thức~\eqref{eq:autoregressive}, có nhiều \textit{chiến lược giải mã}
(\textit{decoding strategy}) để chọn token đầu ra tại mỗi bước. \textit{Giải mã tham lam}
(\textit{greedy}) luôn chọn token có xác suất cao nhất, cho kết quả tất định nhưng dễ rơi vào lặp
hoặc bỏ lỡ lời giải tốt hơn. Các chiến lược lấy mẫu như \textit{top-$k$} (chỉ lấy mẫu trong $k$
token nhiều khả năng nhất) hay \textit{top-$p$} (\textit{nucleus}, lấy mẫu trong tập token nhỏ
nhất có tổng xác suất vượt ngưỡng $p$) đưa thêm tính ngẫu nhiên để tăng đa dạng. Mức ngẫu nhiên
được điều khiển bởi tham số \textit{nhiệt độ} (\textit{temperature}) $T$: giá trị $T$ nhỏ làm phân
phối nhọn hơn (tiệm cận giải mã tham lam), $T$ lớn làm phân phối phẳng hơn. Vì khóa luận cần kết
quả \emph{lặp lại được} (\textit{reproducible}) để so sánh công bằng giữa các cấu hình, mô hình
ngôn ngữ được đặt $T = 0$ trong toàn bộ thực nghiệm.

\subsection{Học trong ngữ cảnh và hiện tượng ảo giác}

Một đặc tính nổi bật của mô hình ngôn ngữ lớn là khả năng \textit{học trong ngữ cảnh}
(\textit{in-context learning}): mô hình có thể thực hiện một tác vụ mới chỉ bằng cách quan sát
vài ví dụ minh họa đặt trong câu lệnh (\textit{prompt}), mà không cần cập nhật tham
số~\cite{zhao2023llmsurvey}. Khả năng này là nền tảng cho việc điều khiển mô hình bằng câu lệnh,
được khai thác xuyên suốt trong khung ARI.

Tuy nhiên, do tham số bị cố định sau huấn luyện và do cơ chế sinh dựa trên xác suất, mô hình ngôn
ngữ lớn thường mắc hiện tượng \textit{ảo giác} (\textit{hallucination}) -- sinh ra thông tin nghe
hợp lý nhưng sai sự thật. Hiện tượng này đặc biệt nghiêm trọng với các câu hỏi đòi hỏi tri thức
thời gian chính xác hoặc tri thức mới phát sinh sau thời điểm huấn luyện. Đây là động lực chính
cho việc kết hợp mô hình ngôn ngữ lớn với nguồn tri thức ngoài có cấu trúc.

\section{Tăng cường tri thức cho mô hình ngôn ngữ lớn}
\label{sec:rag}

Để giảm thiểu ảo giác và bổ sung tri thức cập nhật, một hướng nghiên cứu phổ biến là cung cấp tri
thức ngoài cho mô hình ngôn ngữ lớn. Các phương pháp có thể chia thành hai nhóm: tiêm tri thức
tường minh và tiêm tri thức ngầm~\cite{baek2023kaping}.

\subsection{Sinh tăng cường bằng truy hồi}

\textit{Sinh tăng cường bằng truy hồi} (\textit{Retrieval-Augmented Generation} --
RAG)~\cite{lewis2020rag} là phương pháp tiêm tri thức tường minh tiêu biểu. Ý tưởng là truy hồi
các đoạn văn bản hoặc sự kiện liên quan đến câu hỏi từ một kho tri thức, rồi nối chúng vào câu
lệnh để mô hình tham chiếu khi sinh câu trả lời. RAG giúp làm giàu thông tin cho mô hình, nhưng
hiệu quả phụ thuộc mạnh vào độ chính xác của khâu truy hồi và bị giới hạn bởi độ dài ngữ cảnh: nếu
truy hồi thiếu chính xác sẽ đưa thêm nhiễu, còn nếu lượng tri thức quá lớn thì không thể đưa hết
vào câu lệnh.

\subsection{Câu lệnh tăng cường bằng đồ thị tri thức}

Khi nguồn tri thức là đồ thị tri thức, các phương pháp như KAPING~\cite{baek2023kaping} truy hồi
các bộ ba liên quan đến thực thể trong câu hỏi rồi đưa vào câu lệnh dưới dạng văn bản hóa. Hướng
tiếp cận này là cơ sở cho một trong các phương pháp đối sánh (\textit{baseline}) của khóa luận, gọi
là KG-RAG. Tuy nhiên, với đồ thị tri thức thời gian, một câu hỏi có thể liên quan đến hàng nghìn
sự kiện gắn các mốc thời gian khác nhau, khiến việc nhồi toàn bộ vào câu lệnh là bất khả thi và dễ
gây nhiễu cho quá trình suy luận. Đây chính là điểm mà ARI khắc phục bằng cách \emph{tách bạch}
khâu tương tác tri thức ra khỏi khâu suy luận, thay vì nhồi tri thức trực tiếp vào câu lệnh.

\section{Suy luận có cấu trúc với mô hình ngôn ngữ lớn}
\label{sec:reasoning}

Ngoài việc cung cấp tri thức, một hướng khác nhằm tăng năng lực suy luận của mô hình ngôn ngữ lớn
là dẫn dắt quá trình suy luận theo cấu trúc nhiều bước.

\subsection{Chuỗi suy nghĩ}

Kỹ thuật \textit{chuỗi suy nghĩ} (\textit{Chain-of-Thought} -- CoT)~\cite{wei2022cot} khuyến khích
mô hình sinh ra các bước suy luận trung gian trước khi đưa ra câu trả lời cuối cùng, thay vì trả
lời trực tiếp. Bằng cách ``suy nghĩ thành lời'', mô hình phân rã bài toán phức tạp thành chuỗi
bước nhỏ, cải thiện đáng kể độ chính xác trên các tác vụ đòi hỏi suy luận số học và logic. Tuy
nhiên, CoT thuần túy vẫn dựa trên tri thức tham số của mô hình nên không tránh được ảo giác khi
gặp tri thức thời gian.

\subsection{Suy luận kết hợp hành động}

Khung ReAct~\cite{yao2023react} mở rộng CoT bằng cách đan xen các bước \textit{suy luận}
(\textit{reasoning}) với các bước \textit{hành động} (\textit{action}) tương tác với môi trường
ngoài, chẳng hạn truy vấn một công cụ hoặc một cơ sở tri thức. Sau mỗi hành động, mô hình quan sát
phản hồi từ môi trường rồi tiếp tục suy luận. Cách tiếp cận ``tác tử'' (\textit{agent}) này cho
phép mô hình kết hợp tri thức ngoài một cách có chọn lọc theo từng bước. Khung ARI mà khóa luận kế
thừa cũng vận hành theo mô hình tác tử nhiều bước, nhưng khác biệt ở chỗ mô hình chỉ \emph{lựa
chọn} hành động ở mức trừu tượng, còn việc thực thi hành động trên đồ thị được tách hẳn sang một
mô-đun riêng.

\section{Bộ nhớ và học từ kinh nghiệm}
\label{sec:memory-bg}

\subsection{Bộ nhớ ngoài cho mô hình ngôn ngữ lớn}

Mô hình ngôn ngữ lớn vốn không có bộ nhớ dài hạn, và bộ nhớ ngắn hạn bị giới hạn bởi cửa sổ ngữ
cảnh. Nhiều nghiên cứu đã trang bị cho mô hình một \textit{bộ nhớ ngoài} để lưu lại kinh nghiệm và
tái sử dụng. Reflexion~\cite{shinn2023reflexion} lưu các phản tỉnh về lần thử trước để cải thiện ở
lần thử sau; MemPrompt~\cite{madaan2022memprompt} duy trì một bộ nhớ phản hồi của người dùng để
nhắc lại cho mô hình. Tuy vậy, các phương pháp này phần lớn chỉ \emph{tiếp nhận thụ động} thông
tin lịch sử, chưa chủ động \emph{kiến tạo} tri thức trừu tượng có thể tổng quát hóa cho nhiều bài
toán khác nhau.

\subsection{Thuyết kiến tạo}

Khung ARI lấy cảm hứng từ \textit{thuyết kiến tạo} (\textit{constructivism}) trong giáo
dục~\cite{savery1995constructivism, kirschner2006minimal}. Quan điểm cốt lõi của thuyết kiến tạo
là tri thức không được truyền thụ thụ động từ người dạy sang người học, mà được \emph{chủ động xây
dựng} trong tâm trí người học dựa trên kinh nghiệm sẵn có. Người học tổng hợp tri thức mới bằng
cách liên tục điều chỉnh và tinh lọc khung nhận thức của mình. Áp dụng vào mô hình ngôn ngữ lớn,
ARI hướng tới việc cho mô hình \emph{tự chắt lọc} các chiến lược suy luận trừu tượng từ chính kinh
nghiệm giải toán trước đó -- cả thành công lẫn thất bại -- thay vì chỉ tiếp nhận ví dụ một cách thụ
động. Cơ chế cụ thể hiện thực hóa ý tưởng này sẽ được trình bày ở Mục~\ref{sec:knowledge-agnostic}.

\section{Phân cụm K-means và biểu diễn ngữ nghĩa}
\label{sec:kmeans}

\subsection{Biểu diễn ngữ nghĩa và độ tương đồng cosine}

Để máy tính so sánh ngữ nghĩa giữa các đoạn văn bản, mỗi đoạn được mã hóa thành một vector số thực
trong không gian nhiều chiều, gọi là \textit{biểu diễn nhúng} (\textit{embedding}). Hai đoạn văn
bản có ngữ nghĩa gần nhau sẽ có vector nhúng gần nhau. Độ gần thường được đo bằng \textit{độ tương
đồng cosine} giữa hai vector $\mathbf{u}$ và $\mathbf{v}$:

\begin{equation}
\label{eq:cosine}
\mathrm{cos}(\mathbf{u}, \mathbf{v}) =
\frac{\mathbf{u} \cdot \mathbf{v}}{\lVert \mathbf{u} \rVert \, \lVert \mathbf{v} \rVert}
= \frac{\sum_{i=1}^{d} u_i v_i}{\sqrt{\sum_{i=1}^{d} u_i^2} \, \sqrt{\sum_{i=1}^{d} v_i^2}}.
\end{equation}

Giá trị cosine nằm trong khoảng $[-1, 1]$; càng gần $1$ thì hai vector càng tương đồng. Trong khóa
luận, độ tương đồng cosine được dùng cả ở khâu lọc hành động ứng viên (Mục~\ref{sec:knowledge-based})
lẫn khâu liên kết thực thể (Mục~\ref{sec:vietnamese-adapt}).

\subsection{Thuật toán K-means}

\textit{K-means}~\cite{macqueen1967kmeans} là thuật toán phân cụm không giám sát kinh điển, nhằm
chia một tập điểm dữ liệu thành $k$ cụm sao cho tổng bình phương khoảng cách từ mỗi điểm tới tâm
cụm của nó là nhỏ nhất. Cho tập điểm $\{\mathbf{x}_1, \dots, \mathbf{x}_n\}$, thuật toán tìm $k$
tâm cụm $\{\boldsymbol{\mu}_1, \dots, \boldsymbol{\mu}_k\}$ tối thiểu hóa hàm mục tiêu:

\begin{equation}
\label{eq:kmeans}
J = \sum_{j=1}^{k} \sum_{\mathbf{x}_i \in C_j} \lVert \mathbf{x}_i - \boldsymbol{\mu}_j \rVert^2,
\end{equation}

trong đó $C_j$ là tập điểm thuộc cụm thứ $j$. Thuật toán lặp hai bước cho đến khi hội tụ: (1) gán
mỗi điểm vào cụm có tâm gần nhất; (2) cập nhật lại tâm cụm bằng trung bình các điểm trong cụm.
Trong khung ARI, K-means được dùng để phân loại các trace suy luận lịch sử thành các nhóm câu hỏi
tương tự, từ đó chắt lọc phương pháp luận trừu tượng riêng cho từng nhóm (Mục~\ref{sec:knowledge-agnostic}).

Quá trình lặp của K-means được tóm tắt trong Thuật toán~\ref{alg:kmeans}. Chất lượng phân cụm phụ
thuộc vào cách khởi tạo tâm cụm: kỹ thuật \textit{k-means++} chọn các tâm khởi tạo phân tán nhau
nhằm giảm nguy cơ hội tụ về một cực tiểu cục bộ xấu. Số cụm $k$ là siêu tham số phải chọn trước; số
cụm quá ít khiến mỗi cụm trộn lẫn nhiều dạng câu hỏi khác nhau, còn quá nhiều khiến mỗi cụm có quá
ít mẫu để chắt lọc phương pháp luận tốt. Ảnh hưởng của $k$ tới hiệu năng được khảo sát thực nghiệm
ở Chương~\ref{ch:experiment}.

\begin{algorithm}[ht]
\caption{Thuật toán K-means}
\label{alg:kmeans}
\begin{algorithmic}[1]
\Require Tập điểm $X = \{\mathbf{x}_1, \dots, \mathbf{x}_n\}$, số cụm $k$
\Ensure Phân hoạch $\{C_1, \dots, C_k\}$ và tâm cụm $\{\boldsymbol{\mu}_1, \dots, \boldsymbol{\mu}_k\}$
\State Khởi tạo $k$ tâm cụm (ví dụ bằng k-means++)
\Repeat
  \For{mỗi điểm $\mathbf{x}_i \in X$}
    \State Gán $\mathbf{x}_i$ vào cụm $C_j$ với $j = \arg\min_{l} \lVert \mathbf{x}_i - \boldsymbol{\mu}_l \rVert^2$
  \EndFor
  \For{mỗi cụm $C_j$}
    \State $\boldsymbol{\mu}_j \gets \dfrac{1}{|C_j|} \sum_{\mathbf{x}_i \in C_j} \mathbf{x}_i$
  \EndFor
\Until{các tâm cụm không thay đổi}
\end{algorithmic}
\end{algorithm}

\section{Cross-encoder và bi-encoder}
\label{sec:crossencoder-bg}

Khi cần chấm điểm mức liên quan giữa hai đoạn văn bản (ví dụ một câu hỏi và một ứng viên trả lời),
có hai kiến trúc phổ biến dựa trên mô hình ngôn ngữ tiền huấn luyện như BERT~\cite{devlin2019bert}:
\textit{bi-encoder} và \textit{cross-encoder}~\cite{reimers2019sbert}. Hình~\ref{fig:bi-cross}
minh họa sự khác biệt giữa hai kiến trúc.

\begin{figure}[ht]
  \centering
  \fbox{\parbox{0.92\textwidth}{\centering\vspace{2mm}\footnotesize
  \textbf{Bi-encoder:}\quad $a \rightarrow \mathrm{BERT} \rightarrow \mathbf{e}_a$ \;;\;
  $b \rightarrow \mathrm{BERT} \rightarrow \mathbf{e}_b$ \;;\;
  $\mathrm{score} = \mathrm{cos}(\mathbf{e}_a, \mathbf{e}_b)$\\[3mm]
  \textbf{Cross-encoder:}\quad
  $[\mathtt{CLS}]\,a\,[\mathtt{SEP}]\,b\,[\mathtt{SEP}] \rightarrow \mathrm{BERT}
  \rightarrow \mathrm{Linear} \rightarrow \sigma \rightarrow \mathrm{score}$
  \vspace{2mm}}}
  \caption{So sánh bi-encoder (mã hóa độc lập hai đầu vào rồi đo cosine) và cross-encoder (mã hóa
  đồng thời cặp đầu vào để nắm bắt tương tác ngữ cảnh). \todo{vẽ lại bằng TikZ cho rõ}}
  \label{fig:bi-cross}
\end{figure}

\textbf{Bi-encoder} mã hóa hai đầu vào $a$ và $b$ \emph{một cách độc lập} thành hai vector
$\mathbf{e}_a$, $\mathbf{e}_b$, rồi tính độ liên quan bằng độ tương đồng cosine theo
Công thức~\eqref{eq:cosine}. Ưu điểm là hiệu quả tính toán: các vector có thể được tính trước và
lưu trữ, nên việc so khớp rất nhanh, phù hợp với truy hồi quy mô lớn. Nhược điểm là do mã hóa độc
lập, mô hình không nắm bắt được \emph{tương tác} trực tiếp giữa $a$ và $b$.

\textbf{Cross-encoder} ngược lại nối hai đầu vào thành một chuỗi duy nhất
$[\mathtt{CLS}]\,a\,[\mathtt{SEP}]\,b\,[\mathtt{SEP}]$ rồi đưa qua mô hình BERT, sau đó dùng một
lớp tuyến tính kèm hàm sigmoid $\sigma(\cdot)$ để xuất ra điểm liên quan trong khoảng $(0, 1)$:

\begin{equation}
\label{eq:crossencoder}
\mathrm{score}(a, b) = \sigma\!\left(\mathbf{w}^{\top}\,
\mathrm{BERT}_{\mathtt{CLS}}\big([\mathtt{CLS}]\,a\,[\mathtt{SEP}]\,b\,[\mathtt{SEP}]\big) + c\right),
\end{equation}

trong đó $\mathbf{w}$ và $c$ là tham số của lớp tuyến tính, còn
$\mathrm{BERT}_{\mathtt{CLS}}(\cdot)$ là vector biểu diễn tại vị trí $[\mathtt{CLS}]$. Vì hai đầu
vào được mã hóa \emph{đồng thời}, cross-encoder nắm bắt được tương tác ngữ cảnh giữa chúng, cho độ
chính xác cao hơn bi-encoder, đổi lại chi phí tính toán lớn hơn do không thể tính trước. Trong khóa
luận, đặc tính \emph{nhận biết ngữ cảnh} này là cơ sở cho bộ lọc hành động đề xuất ở
Mục~\ref{sec:crossencoder}: thay vì chỉ so khớp hành động với câu hỏi gốc bằng cosine, cross-encoder
cho phép chấm điểm hành động \emph{trong bối cảnh lịch sử suy luận hiện tại}.

\subsection{Huấn luyện và đánh đổi}

Bộ chấm điểm cross-encoder được huấn luyện như một bộ phân loại nhị phân: mỗi mẫu là một cặp
$(a, b)$ kèm nhãn $y \in \{0, 1\}$ cho biết $a$ có liên quan tới $b$ hay không, và mô hình tối
thiểu hóa hàm mất mát \textit{entropy chéo nhị phân} (\textit{binary cross-entropy}):

\begin{equation}
\label{eq:bce}
\mathcal{L} = -\frac{1}{N} \sum_{i=1}^{N}
\Big[\, y_i \log s_i + (1 - y_i) \log (1 - s_i) \,\Big],
\end{equation}

trong đó $s_i = \mathrm{score}(a_i, b_i)$ là điểm dự đoán theo Công thức~\eqref{eq:crossencoder} và
$N$ là số mẫu huấn luyện. Chất lượng mô hình phụ thuộc nhiều vào cách lấy \textit{mẫu âm}
(\textit{negative sampling}) -- tức chọn các cặp không liên quan đủ ``khó'' để mô hình học được
ranh giới phân biệt tinh tế.

Bảng~\ref{tab:bicross} tóm tắt đánh đổi giữa hai kiến trúc. Trong thực tế, một mô thức phổ biến là
dùng bi-encoder để \emph{truy hồi nhanh} một tập ứng viên rộng, rồi dùng cross-encoder để
\emph{xếp hạng lại} (\textit{re-rank}) tập ứng viên nhỏ đó cho chính xác. Đây cũng chính là mô thức
mà bộ lọc hành động đề xuất trong khóa luận áp dụng (Mục~\ref{sec:crossencoder}).

\begin{table}[ht]
\centering
\caption{Đánh đổi giữa bi-encoder và cross-encoder.}
\label{tab:bicross}
\begin{tabular}{L{3.5cm} C{2.8cm} C{2.8cm} C{2.8cm}}
\toprule
\textbf{Kiến trúc} & \textbf{Mã hóa} & \textbf{Tính trước} & \textbf{Độ chính xác} \\
\midrule
Bi-encoder & Độc lập & Được & Thấp hơn \\
Cross-encoder & Đồng thời & Không & Cao hơn \\
\bottomrule
\end{tabular}
\end{table}

\section{Nhận dạng thực thể có tên cho tiếng Việt}
\label{sec:ner-bg}

\textit{Nhận dạng thực thể có tên} (\textit{Named Entity Recognition} -- NER) là bài toán xác định
và phân loại các cụm từ chỉ thực thể (người, tổ chức, địa điểm, \dots) trong văn bản. Trong hệ
thống TKGQA, NER đóng vai trò bước đầu của \textit{liên kết thực thể} (\textit{entity linking}):
trích xuất các cụm chỉ thực thể trong câu hỏi để ánh xạ về đúng nút trong đồ thị tri thức.

Với tiếng Việt, các mô hình ngôn ngữ tiền huấn luyện như PhoBERT~\cite{nguyen2020phobert} hay các
biến thể dựa trên ELECTRA~\cite{clark2020electra} đã đạt chất lượng cao trên nhiều tác vụ, trong đó
có NER. Khóa luận sử dụng một mô hình NER tiếng Việt tiền huấn luyện sẵn có để trích xuất thực thể
từ câu hỏi mở của người dùng, sau đó kết hợp với tìm kiếm tương đồng trong cơ sở dữ liệu vector để
khớp về mã thực thể trong đồ thị; chi tiết được trình bày ở Mục~\ref{sec:vietnamese-adapt}.

\section{Tổng kết chương}

Chương này đã trình bày các nền tảng lý thuyết của khóa luận: từ đồ thị tri thức thời gian và phát
biểu hình thức của bài toán TKGQA, qua mô hình ngôn ngữ lớn cùng các kỹ thuật tăng cường tri thức,
suy luận có cấu trúc và bộ nhớ -- trong đó thuyết kiến tạo là nền tảng tư tưởng của ARI, đến các
công cụ kỹ thuật như phân cụm K-means, biểu diễn ngữ nghĩa, cross-encoder và NER tiếng Việt. Những
khái niệm này sẽ được vận dụng trực tiếp khi mô tả phương pháp đề xuất ở Chương~\ref{ch:method}.
Trước đó, Chương~\ref{ch:related} điểm lại các công trình liên quan và định vị đóng góp của khóa
luận trong bức tranh nghiên cứu chung.

% ===================== Bảng ký hiệu ===================== %
\begin{table}[ht]
\centering
\caption{Các ký hiệu toán học dùng xuyên suốt khóa luận.}
\label{tab:notation}
\begin{tabular}{c L{9.5cm}}
\toprule
\textbf{Ký hiệu} & \textbf{Ý nghĩa} \\
\midrule
$\mathcal{K}$ & Đồ thị tri thức thời gian $(\mathcal{E}, \mathcal{R}, \mathcal{T}, \mathcal{F})$ \\
$\mathcal{E}, \mathcal{R}, \mathcal{T}$ & Tập thực thể, quan hệ, mốc thời gian \\
$\mathcal{F}$ & Tập sự kiện thời gian (bộ bốn) \\
$(s, r, o, \tau)$ & Một sự kiện: chủ thể, quan hệ, đối tượng, thời gian \\
$q$ & Câu hỏi đầu vào \\
$a$ & Đáp án (thực thể hoặc mốc thời gian) \\
$e_h$ & Thực thể chủ thể (\textit{head}) của câu hỏi \\
$G_{e_h}$ & Đồ thị con một bước của $e_h$ \\
$P_t, P'_t$ & Tập hành động ứng viên (trước và sau khi lọc) tại bước $t$ \\
$H_q$ & Tập lịch sử suy luận của câu hỏi $q$ \\
$C_j, \boldsymbol{\mu}_j$ & Cụm thứ $j$ và tâm của nó \\
$M_{C^*}$ & Phương pháp luận trừu tượng của cụm phù hợp nhất $C^*$ \\
$\mathrm{cos}(\cdot, \cdot)$ & Độ tương đồng cosine \\
\bottomrule
\end{tabular}
\end{table}
